\section{Contributions}
\label{sec:contributions}


This dissertation, \emph{Securing High-Stakes Software Against Untrusted
Development and Untrusted Use}, studies how to protect software that guards
critical assets when neither the code that builds it nor the parties that use it
can be fully trusted. As the motivating examples show, harmful behavior can be
introduced while software is written, by AI coding tools that fabricate or follow
untrusted content, and it can be triggered after software is deployed, by
interactions its designers never anticipated. In both settings, the harmful
behavior is induced by factors that cannot be fully enumerated or validated in
advance.

We organize the work around these two settings. The first half of the
dissertation concerns untrusted development, where the goal is to prevent or
detect harmful code before it is trusted or deployed. The second half concerns
untrusted use, where the goal is to protect already-deployed software against
adversarial interaction at runtime. We instantiate each setting in a concrete
domain: LLM-assisted coding systems for development, and smart contracts for
deployment. \revise{Recall the high-stakes software lifecycle of
Figure~\ref{fig:workflow}; Figure~\ref{fig:contributions} revisits that same
lifecycle, now mapping each contribution onto the stage it defends.}

\begin{figure}[t]
  \centering
  \includegraphics[width=\textwidth]{figures/workflow2.pdf}
  \caption{\revise{Overview of the contributions of this dissertation, organized
  along the lifecycle of high-stakes software. The development side (untrusted
  development) covers defending coding agents against search- and
  documentation-level poisoning (Chapter~2) and auditing malicious endpoints in
  production LLMs (Chapter~3); the deployment side (untrusted use) covers
  control-flow integrity (Chapter~4) and runtime invariants (Chapter~5) for
  deployed smart contracts.}}
  \label{fig:contributions}
\end{figure}

This thesis makes the following contributions.
\begin{itemize}
    % \item \textbf{A unified defense framing for high-stakes software
    % under untrusted development and use.}
    % We formulate a common perspective for software whose behavior may be shaped
    % by AI-generated code, weakly trusted external content, opaque models, or
    % adversarial interaction, and we identify defense opportunities both before
    % and after deployment.

    \item \textbf{Defending coding agents against search- and
    documentation-level poisoning (Chapter~\revise{2}).}
    We study how poisoned external content retrieved by coding agents propagates
    into harmful coding outcomes, and we develop defenses that reduce this risk
    while preserving normal agent utility. This is the planned final work of the
    dissertation.

    \item \textbf{Auditing malicious endpoints in production LLMs (Chapter~\revise{3}).}
    We audit production language models as they generate code under ordinary
    developer-style prompts, and show that they recommend malicious resources,
    such as scam endpoints, at non-trivial rates. We build an automated framework
    that constructs benign-looking coding requests at scale and measures this
    risk~\cite{chen2025scam2prompt}.

    \item \textbf{Control-flow integrity for smart contracts (Chapter~\revise{4}).}
    We propose algorithms that enforce control-flow integrity across contract
    interactions with low runtime overhead, ensuring that contract functions are
    invoked only in legitimate patterns~\cite{chen2025secure}.

    \item \textbf{Runtime invariants for deployed smart contracts (Chapter~\revise{5}).}
    We propose techniques that learn behavioral invariants from historical
    transaction traces and enforce them as runtime guards, stopping transactions
    that depart from normal use~\cite{chen2024demystifying}.
\end{itemize}
